\documentclass[10pt,a4paper]{amsart}
\usepackage[margin=19mm]{geometry}
\usepackage{amsmath,amssymb,mathtools}
\usepackage[scr=rsfs,cal=euler]{mathalfa}
\usepackage{xfrac}
\usepackage[unicode,hidelinks,bookmarks=false]{hyperref}
\newcommand{\ie}{i.e.\ }
\newcommand{\cf}{cf.\ }
\newcommand{\resp}{respectively\ }
\newcommand{\Subsection}{\subsection}
\providecommand{\nobreakdash}{\nobreak\hbox{-}}
\setlength{\emergencystretch}{2em}
\multlinegap0pt
% ---------------------------------------------------------------------------
% Editorial AMS wrapper prelude. The theorem environments and the mathematical macros below
% are transcribed from the paper's own preamble (cached-originals/source0.tex lines 24--120); the AMS
% amsart class, the named format packages of the pinned TeX Live runtime and this editorial formatting
% are not part of the author's text. No mathematics is changed.
% ---------------------------------------------------------------------------
\newcommand{\abs}[1]{|#1|^2}
\newcommand{\Abs}[1]{\Big|#1\Big|^2}
\newcommand{\bas}[1]{|#1|^2}
\newcommand{\inp}[1]{\langle #1\rangle}



\def\Xint#1{\mathchoice
	{\XXint\displaystyle\textstyle{#1}}%
	{\XXint\textstyle\scriptstyle{#1}}%
	{\XXint\scriptstyle\scriptscriptstyle{#1}}%
	{\XXint\scriptscriptstyle\scriptscriptstyle{#1}}%
	\!\int}
\def\XXint#1#2#3{{\setbox0=\hbox{$#1{#2#3}{\int}$ }
		\vcenter{\hbox{$#2#3$ }}\kern-.6\wd0}}
\def\ddashint{\Xint=}
\def\dashint{\Xint-}
\newtheorem{claim}{Claim}
\newtheorem{lem}{Lemma}[section]
\newtheorem{prop}{Proposition}[section]
\newtheorem{ques}{Question}
\newtheorem{example}{Example}

\newtheorem{defn}{Definition}[section]
\newtheorem{corr}{Corollary}[section]
\newtheorem{conj}{Conjecture}
\newtheorem{remark}{Remark}[section]
\newtheorem{theorem}{Theorem}[section]

\newcommand{\medent}{\bigskip\noindent $\blacktriangleright$~}
\allowdisplaybreaks
\newcommand{\del}{\partial}
\newcommand{\dbar}{\bar\partial}
\newcommand{\eps}{\varepsilon}
\newcommand{\ddbar}{\sqrt{-1}\partial\bar\partial}
\newcommand{\innpro}[1]{\langle#1\rangle}
\newcommand{\bk}[1]{\Big(#1\Big)}
\newcommand{\ppm}[1]{\begin{pmatrix}#1\end{pmatrix}}
\newcommand{\xk}[1]{\big(#1\big)}
\newcommand{\C}{\mathbb C}
\newcommand{\R}{\mathbb R}

\newcommand{\detla}{\delta}
\newcommand{\Detla}{\Delta}
\newcommand{\lapl}{\Delta}
\newcommand{\minus}{\backslash}
\newcommand{\ssubset}{\subset\joinrel\subset}
\newcommand{\ssupset}{\supset\joinrel\supset}

\makeindex
\DeclareMathOperator{\diam}{diam}
\DeclareMathOperator{\ric}{Ric}
\DeclareMathOperator{\wg}{\wedge}
\DeclareMathOperator{\im}{\sqrt{-1}}
\DeclareMathOperator{\isom}{\simeq}
\DeclareMathOperator{\Span}{Span}

\DeclareMathOperator{\osc}{osc}
\DeclareMathOperator{\vol}{Vol}
\DeclareMathOperator{\re}{\text{Re}}

\DeclareMathOperator{\dimm}{dim}
\DeclareMathOperator{\kernel}{Ker}
\DeclareMathOperator{\psh}{PSH}
\DeclareMathOperator{\tr}{tr}
\DeclareMathOperator{\capa}{cap}
\DeclareMathOperator{\spec}{Spec}

\renewcommand{\theequation}{\thesection.\arabic{equation}}
\numberwithin{equation}{section}



\begin{document}
\title{Boundary regularity of optimal transport maps on convex domains}
\author{Tristan C. Collins}
\email{tristanc@math.toronto.edu}
\author{Freid Tong}
\email{freid.tong@utoronto.ca}
\address{Department of Mathematics, University of Toronto, 40 St. George St., Toronto, Ontario, Canada}
\date{Source: arXiv:2507.05395v2; distributed under CC BY 4.0}
\maketitle
\noindent\emph{Source, version and licence.} \emph{Boundary regularity of optimal transport maps on convex domains}, by Tristan C. Collins and Freid Tong. The source of this editable unit is the e-print of \textbf{version v2} of arXiv:2507.05395, https://arxiv.org/abs/2507.05395v2, whose material is distributed under the Creative Commons Attribution 4.0 International licence, http://creativecommons.org/licenses/by/4.0/, as recorded in the arXiv licence metadata for this paper and shown on that abstract page. The whole of the body below is version v2, the newest version of the paper. Full rights notice: licenses/arxiv-2507.05395-CC-BY-4.0.txt.\par\smallskip
\noindent\textit{Editorial scope.} Counted unit: original Proposition 7.1 of the article, the statement carried by the source environment \texttt{prop} with source label \texttt{prop: VeryGenPog}, cached-originals/source0.tex lines 1260--1294, printed as \textbf{Proposition 7.1} exactly as in the article. It is a Pogorelov-type estimate along flat directions: under the four structural conditions on the domains, the densities and the boundary data, the second derivative $u_{11}$ of the convex solution of the Monge--Amp\`ere problem is controlled by $|u-\ell-h|$ on the section $S_h(u,0)$, uniformly in the data. It is delivered together with the author's own complete proof as the single primary proof body, source0.tex lines 1295--1432 (138 lines): the \texttt{proof} environment that immediately follows the statement, that runs the Pogorelov maximum-principle argument for the quantity $\Phi=\log u_{11}+\log|u|+\frac12 u_1^2$, treats the interior critical point case and the boundary critical point case with the boundary data and the shearing transformation, and closes with the desired estimate. No proof marker is added and no mathematics is written, repaired or re-derived; the author's notation and the printed slips of the source are kept literal. Necessary complete context is retained uncounted: source0.tex lines 279--356, the article's own subsection ``Sections, centered sections, and extrinsic balls'', which introduces the sections $S_h(u,x)$, the centered sections and the extrinsic balls that the counted statement and its proof use, kept with its original subsection numbering. The article's own numbering is reproduced by explicit counter settings: this article gives each theorem-like environment its own section-scoped counter (\texttt{\textbackslash newtheorem\{prop\}\{Proposition\}[section]}, and likewise for theorem, lem, defn, corr and remark), and every one of those counters together with the section-scoped equation counter is set before each retained passage, so the counted statement prints as Proposition 7.1, the retained displays keep their original tags (7.1), (7.2), (7.3) and the following ones, and the retained subsection prints with its original number. Every \texttt{\textbackslash ref} and \texttt{\textbackslash eqref} inside every retained passage resolves inside this delivered file, and the seven citations that occur in the retained passages are delivered as the author's own bibliography entries transcribed verbatim from the article's own \texttt{thebibliography} environment, keeping the author's own printed numbers [3], [4], [5], [9], [18], [26] and [35]. The e-print ships no artwork file and no \texttt{\textbackslash includegraphics} command occurs in any retained passage; the article's figures and diagrams lie outside every retained passage, so no figure environment and no graphics-inclusion command occurs anywhere in this delivered file and no artwork is redistributed. The source's own class is AMS \texttt{amsart} (\texttt{\textbackslash documentclass[12pt,reqno]\{amsart\}}); the e-print ships no class or style file, so no publisher class is shipped, used or redistributed, and the wrapper is the same AMS \texttt{amsart} class with the paper's own macros and theorem declarations transcribed from its preamble. Every declared body of this unit is an exact, unmodified span of source0.tex: no body is a bounded passage comparison, \texttt{omit} was not used anywhere, and consequently no fidelity receipt is asserted in delivery-evidence.json. The complete concatenated original source is bundled as sources/181172/source0.tex. Omitted from this editable unit and therefore absent from it are source0.tex lines 1--278; 357--1259; 1433--1664; 1671--1679; 1682--1698; 1701--1714; 1718--1733; 1736--1736, that is the article's own preamble and title block, the introduction with its main theorems, the monotonicity-formula section, the homogeneity-of-blow-ups section, the global regularity sections, the roundness section, the remaining part of the flat-directions section including the mixed-homogeneity shape theorem and the deferred volume lemma, and the remaining bibliography entries, all of which are present in the bundled source but not delivered here. The AMS \texttt{amsart} wrapper, the named format packages of the pinned TeX Live runtime and this editorial source, licence and scope paragraph are editorial formatting only.\par\smallskip
\setcounter{section}{2}
\setcounter{equation}{3}
\setcounter{claim}{0}
\setcounter{conj}{0}
\setcounter{corr}{0}
\setcounter{defn}{4}
\setcounter{example}{0}
\setcounter{lem}{0}
\setcounter{prop}{2}
\setcounter{ques}{0}
\setcounter{remark}{1}
\setcounter{theorem}{0}
\subsection{Sections, centered sections, and extrinsic balls}
Let $((\Omega, g(x)dx), (\Omega', g'(y)dy), u, v)$ be an optimal transport map. Following Caffarelli \cite{Caffarelli}, for any $x_0\in \overline\Omega$, $\mathtt{y}_0\in\mathbb{R}^n$ and $h>0$, we define the {\bf section of $u$ at $x_0$ of height $h$ in direction $\mathtt{y}_0$} to be the convex set
\[S_h(u, x_0, \mathtt{y}_0) := \{x\in\Omega: u(x)-u(x_0)-\mathtt{y}_0(x-x_0)\leq h\}.\]
For simplicity, we shall denote
\[
S_{h}(u,x_0):= S_{h}(u,x_0, \nabla u(x_0))
\]
We also define a {\bf centered section} of $u$ at $x_0\in\overline\Omega$ to be
\[S^c_h(u, x_0) := \{x\in \mathbb R^n:\bar u(x)-\bar u(x_0)-p\cdot (x-x_0)\leq h\}\]
where recall that $\bar u$ is the minimal convex extension of $u$, and $p\in \mathbb R^n$ is chosen so that $S^c_h(u, x_0)$ has center of mass $x_0$. By \cite{Caffarelli2}, if the graph of $\bar u$ contains no line, then centered sections always exist. 

Given a centered section $S = S^c_h(u, x_0)$, one can define the normalized section by
\[\tilde S = A(S)\]
where $A$ is the (unique) symmetric positive definite matrix such that $A(E) = B_1(0)$ where $E$ is the John ellipsoid of $S$. We call $A$ the {\bf normalization matrix} of $S$, and it satisfies $(\det A)|S|\sim 1$. We also define the normalized potential $\tilde u:\tilde S\to \mathbb R$ by
\[\tilde u(\tilde x) = \frac{(\bar u-l)(A^{-1}\tilde x)}{h}.\]
where $l(x) = u(x_0)+p\cdot (x-x_0)$. It follows that $\tilde S = \{\tilde u<1\}$, and the pair $(\tilde S, \tilde u)$ is a normalized pair. Moreover, $\nabla \tilde u:\tilde\Omega\to\tilde\Omega'$ where $\tilde\Omega = A(\Omega)$ and $\tilde\Omega' = hA^{-1}(\Omega')$.  

Furthermore, the following properties were established for doubling measures by Caffarelli \cite{Caffarelli2, Caffarelli3} and Jhaveri-Savin \cite{Jhaveri-Savin}. 


\begin{prop}\label{prop: properties-of-sections}\cite[Section 3]{Jhaveri-Savin}
There exist constant $c>0$ depends only on dimension, and $C\geq 1$ depending additionally on the doubling constant of $g(x)$ and $g'(y)$ such that
\begin{enumerate}
    \item $B_{C^{-1}}\subset \nabla\tilde u(\tilde S)\subset B_C$
    \[\implies C^{-1}h(S_h^c(u, x_0)-x_0)^{\circ}\subset \nabla (u-l)(S_h^c(u, x_0))\subset Ch(S_h^c(u, x_0)-x_0)^{\circ}.\]
    \item $C^{-1}h^n\leq|S_h^c(u, x_0)||\nabla u(S_h^c(u, x_0))|\leq Ch^n$. 
    \item $S_{ch}^c(u, x_0)\cap \overline\Omega\subset S_h(x_0)\subset S_{Ch}^c(u, x_0)\cap \overline\Omega$. 
    \item $S_{C^{-1}h}(v, \nabla u(x_0))\subset \nabla u(S_h(u, x_0))\subset S_{Ch}(v, \nabla u(x_0))$. 
\end{enumerate}
\end{prop}

In \cite{Caffarelli3, Jhaveri-Savin}, it was also shown that if the measures are doubling, then the centered sections satisfy the following engulfing property. 
\begin{prop}\cite[Lemma 3.3]{Jhaveri-Savin}\label{prop: engulfing}
For any pair of constants $0\leq \underline{t}\leq \overline t\leq 1$, there exists a constant $0<s\leq 1$ such that
\[S_{sh}^c(u, x')\subset \overline tS_h^c(u, x)\]
for all $x\in \overline\Omega$ and $x'\in \underline{t}S_h^c(u, x)\cap \overline\Omega$. Moreover, the constant $s$ depends only on $\underline{t}, \overline t, n$ and the doubling constant of $g$ and $g'$.
\end{prop}

This engulfing property leads to the following $C^{1, \delta}$ estimate \cite{Caffarelli2, Jhaveri-Savin}. (see also \cite{Savin-Yu})    

\begin{prop}\label{prop: C1-delta-estimate}
        Given an optimal transport map $((\Omega, g(x)dx), (\Omega', g'(y)dy), u, v)$ with $0\in\partial\Omega\cap \partial\Omega'$ and $u(0) = |\nabla u|(0) = 0$. Assume that it satisfies
        \begin{enumerate}
            \item $g$ and $g'$ has doubling constant bounded by $K$.
            \item $S_1^c(u, 0)\subset B_K(0)$ and $S_1^c(v, 0)\subset B_K(0)$. 
        \end{enumerate}
        Then there exist $\delta>0$ depending only on $K$ and $n$ such that
        \[\|\bar u\|_{C^{1, \delta}(B_R)}+\|\bar v\|_{C^{1, \delta}(B_R)}\leq C \]
        for $C$ depending only on $n, K, $ and $R$. 
    \end{prop}
    \begin{proof}
        The proof follows from \cite[Proposition 2.5, Remark 2.1]{Savin-Yu} using Proposition~\ref{prop: engulfing}. 
    \end{proof}

Moreover, we will also define the following sets 
\[D_r(u, x_0,y_0) := \{x\in\Omega: (x-x_0)\cdot (\nabla u(x)-y_0)\leq r^2\}\]
which we call {\bf extrinsic balls of radius $r$ centered at $(x_0,y_0)$}.  For simplicity, when $y_0=\nabla u(x_0) $, we shall write $D_r(u, x_0,\nabla u(x_0))=: D_r(u,x_0)$.  These extrinsic balls were also used in \cite{Yuan}. 
The following proposition shows that up to a factor of 2, the extrinsic ball and the sections have the same shape. 
\begin{prop}\label{prop: comparison between sections and extrinsic ball}
Assume that $u$ is strictly convex, and $\underline{x}$ is chosen to that $\nabla u(\underline{x}) = \mathtt{y}_0$ then 
    \[\frac{1}{2}S_{r^2}(u, \underline{x})\subset D_r(u, x_0,\mathtt{y}_0)\subset S_{r^2}(u, x_0, \mathtt{y}_0).\]
    In particular, we have
    \[
    \frac{1}{2}S_{r^2}(u, x_0)\subset D_r(u, x_0)\subset S_{r^2}(u, x_0).
    \]
\end{prop}
\begin{proof}
    Without loss of generality, we can assume $x_0 = 0$, $u(0) =0$, and $\nabla u(0) = 0$ by a shift and subtracting a linear function.  Let $\tilde{u} =u(x)- \langle \mathtt{y}_0,x\rangle$. It follows by convexity that 
    \[
        \tilde{u}(x)=\tilde{u}(x)-\tilde{u}(0) \leq x\cdot \nabla\tilde{u}(x) \leq \tilde{u}(2x)- \tilde{u}(x)
        \] 
        This yields the containment $D_r(u, x_0,\mathtt{y}_0)\subset S_{r^2}(u, x_0, \mathtt{y}_0)$.  Furthermore, if $\mathtt{y}_0=0$, then $\tilde{u}(x)=u(x) \geq 0$ and the second claim follows.  It only remains to prove the containment $\frac{1}{2}S_{r^2}(u, \underline{x})\subset D_r(u, x_0,\mathtt{y}_0)$.  For this, suppose $2x\in S_{r^2}(u,\underline{x})$.  Then we have
        \[
        r^2 \geq \tilde{u}(2x) - \tilde{u}(\underline{x}) \geq \nabla \tilde{u}(x) \cdot x + \tilde{u}(x) - \tilde{u}(\underline{x})
        \]
        On the other hand, by definition of $\underline{x}$ we have $\tilde{u}(x) - \tilde{u}(\underline{x}) \geq 0$, and the result follows.
\end{proof}


\setcounter{section}{7}
\setcounter{equation}{0}
\setcounter{claim}{0}
\setcounter{conj}{0}
\setcounter{corr}{0}
\setcounter{defn}{0}
\setcounter{example}{0}
\setcounter{lem}{0}
\setcounter{prop}{0}
\setcounter{ques}{1}
\setcounter{remark}{0}
\setcounter{theorem}{0}
\begin{prop}\label{prop: VeryGenPog}
Let $\Omega\subset \mathbb{R}^{n}_{x}$ and $ \Omega' \subset \mathbb{R}_y^n$.  Suppose that the following structural conditions hold: 
\begin{itemize}
    \item[(i)] There is a smooth convex function $\phi_\Omega(x) =\phi_\Omega(x_2,\ldots, x_{n-1})$ such that $\phi_\Omega(0)=0$ and
    \[
    \Omega \cap B_1(0) = \{ x\in B_1(0) : x_n \geq \phi_\Omega(x_2,\ldots x_{n-1})\}
    \]
    \item[(ii)] There is a smooth convex function $\phi_{\Omega'}(y)=\phi_{\Omega'}(y_1,\ldots, y_{n-1})$ such that $\phi_{\Omega'}(0)=0$ and
    \[
    \Omega'\cap B_1(0) =  \{ y\in B_1(0) : y_n \geq \phi_{\Omega'}(y_1,\ldots y_{n-1})\}
    \]
    \item[(iii)] There is a smooth, positive function $g(x)$ defined on $\overline{\Omega\cap B_1(0)} $ and such that $g(x)=g(x_2,\ldots,x_n)$ is independent of $x_1$.
    \item[(iv)]There is a smooth, positive function $g'(y)$ defined on $\overline{\Omega'\cap B_1(0)} $ and such that $g'(y)$ is log-concave and satisfies
    \[
    \sum_{p=1}^{n}\frac{\del g'}{\del y_p}y_p\leq \kappa g'(y)
    \]
    for some $\kappa >0$ and for all $y\in \overline{\Omega'\cap B_1(0)}$.
\end{itemize}
Suppose that $u$ is a convex function satisfying $0 = u(0) = |\nabla u|(0)$ and 
\[
\begin{aligned}
g'(\nabla u)\det D^2u &= g(x),& \quad \text{ in } \Omega \cap B_1(0)\\
\frac{\del u}{\del x_n} &\geq \phi_{\Omega'}(\frac{\del u}{\del x_1}, \ldots, \frac{\del u}{\del x_{n-1}}),& \quad \text{ in } \Omega \cap B_1(0)\\
\frac{\del u}{\del x_n} &= \phi_{\Omega'}(\frac{\del u}{\del x_1}, \ldots, \frac{\del u}{\del x_{n-1}}),& \quad \text{ on } \del\Omega \cap B_1(0)
\end{aligned}
\]
Let $\ell$ be the tangent plane to $u$ at the origin, $h>0$ and suppose that
\[
S_{h}(u,0) = \{u< \ell+h\} \Subset B_1(0)\cap \Omega.
\]
Then we have a bound
\[
u_{11}|u-\ell-h| \leq C(n,\kappa, \|u_{1}-\ell_{1}\|_{L^{\infty}(S_{h}(0))}). 
\]
\end{prop}

\setcounter{section}{7}
\setcounter{equation}{0}
\setcounter{claim}{0}
\setcounter{conj}{0}
\setcounter{corr}{0}
\setcounter{defn}{0}
\setcounter{example}{0}
\setcounter{lem}{0}
\setcounter{prop}{1}
\setcounter{ques}{1}
\setcounter{remark}{0}
\setcounter{theorem}{0}
\begin{proof}
    By subtracting a linear function from $u$, we may assume that $u(0)= |\nabla u(0)|=0$.  Let $S_h= S_h(u,0)$ be the a section of height $h$ centered at $0$, and replace $u$ with $u-h$.  The argument is a Pogorelov type argument, making essential use of the boundary data.  Consider the quantity
\[
\Phi = \log u_{11} + \log|u| + \frac{1}{2}u_1^2
\]
    Let $x_* \in S_h$ be the point where $\Phi$ achieves its maximum.  Either $x_*$ is an interior point, or $x_*$ lies on $\del\Omega \cap S_h\cap B_1$. 

Suppose first that $x_*$ is an interior point. We may perform a shearing transformation preserving the $x_1$ direction, followed by a rotation in the $x_2,\ldots,x_n$ directions in order to diagonalize $D^2u(x_*)$.  We note that since the $x_1$-direction is preserved, the structural conditions $(i)-(iv)$ are preserved under this affine change of coordinates. Differentiating the equation in the $x_1$-direction and using that $D^2u$ is diagonal at $x_*$  we compute
\begin{equation}\label{eq: LinearizedEquation}
\sum_p \frac{u_{1pp}}{u_{pp}}  + \frac{1}{g'}\frac{\del g'}{\del y_1}u_{11} = 0
\end{equation}
\begin{equation}\label{eq: concavityEquation}
\sum_p \frac{u_{11pp}}{u_{pp}} + \frac{1}{g'}\frac{\del g'}{\del y_p}u_{11p} = \sum_{p,q}\frac{u_{1pq}^2} {u_{pp}u_{qq}}-u_{11}^2\frac{\del^2 \log g'}{\del y_1^2}
\end{equation}

Differentiating $\Phi$ at $x_*$ we have
\begin{equation}\label{eq: firstDerivTestPhi}
\Phi_{p} = \frac{u_{11p}}{u_{11}} + \frac{u_p}{u} + u_{1p}u_1=0
\end{equation}
\begin{equation}\label{eq: secondDerivTestPhi}
 \Phi_{pp} = \frac{1}{u_{11}}u_{11pp}-\frac{(u_{11p})^2}{u_{11}^2} + \frac{u_{pp}}{u} - \frac{u_p^2}{u^2} + u_{1pp}u_1 + u_{1p}^2 \leq0
\end{equation}
Multiplying~\eqref{eq: secondDerivTestPhi} by $u^{pp}= \frac{1}{u_{pp}}$ and summing over $p$ yields
\[
0 \geq \sum_{1\leq p\leq n} \frac{u_{11pp}}{u_{11}u_{pp}}-\sum_{p=1}^{n}\frac{(u_{11p})^2}{u_{11}^2u_{pp}} + \frac{n}{u} - \sum_{p=1}^{n}\frac{u_p^2}{u^2u_{pp}} + u_{1}\sum_{p=1}^{n}\frac{u_{1pp}}{u_{pp}} + u_{11} 
\]
We use~\eqref{eq: firstDerivTestPhi} to replace the fourth term on the right hand side for $p\ne 1$ to get
\[
0 \geq \sum_{1\leq p\leq n} \frac{u_{11pp}}{u_{11}u_{pp}}-\sum_{p=1}^{n}\frac{(u_{11p})^2}{u_{11}^2u_{pp}} + \frac{n}{u} - \sum_{p=2}^{n}\frac{u_{11p}^2}{u_{11}^2u_{pp}}-\frac{u_1^2}{u^2u_{11}} + u_{1}\sum_{p=1}^{n}\frac{u_{1pp}}{u_{pp}} + u_{11}
\]
We now use~\eqref{eq: concavityEquation} to replace the fourth order terms
\begin{equation}\label{eq: PogorelovInequality1}
\begin{aligned}
  0 &\geq \sum_{1 \leq p,q\leq n}\frac{u_{1pq}^2} {u_{pp}u_{qq}u_{11}}-u_{11}\frac{\del^2 \log g'}{\del y_1^2}+\frac{1}{u_{11}}\frac{\del^2\log g}{\del x_1^2} -\sum_{p=1}^{n}\frac{1}{g'u_{11}}\frac{\del g'}{\del y_p}u_{11p} \\
  &\quad -\sum_{p=1}^{n}\frac{(u_{11p})^2}{u_{11}^2u_{pp}} + \frac{n}{u} - \sum_{p=2}^{n}\frac{u_{11p}^2}{u_{11}^2u_{pp}}-\frac{u_1^2}{u^2u_{11}} + u_{1}\sum_{p=1}^{n}\frac{u_{1pp}}{u_{pp}} + u_{11}\\
  &\geq \sum_{2 \leq p,q\leq n}\frac{u_{1pq}^2} {u_{pp}u_{qq}u_{11}}-u_{11}\frac{\del^2 \log g'}{\del y_1^2} -\sum_{p=1}^{n}\frac{1}{g'u_{11}}\frac{\del g'}{\del y_p}u_{11p} \\
  &\quad + \frac{n}{u} -\frac{u_1^2}{u^2u_{11}} + u_{1}\sum_{p=1}^{n}\frac{u_{1pp}}{u_{pp}} + u_{11}
\end{aligned}
\end{equation}
We now use~\eqref{eq: firstDerivTestPhi} to obtain
\[
\begin{aligned}
-\sum_{p=1}^{n}\frac{1}{g'u_{11}}\frac{\del g'}{\del y_p}u_{11p} &= \sum_{p=1}^{n}\frac{1}{g'}\frac{\del g'}{\del y_p}\left(\frac{u_{p}}{u}+u_{1p}u_1\right)\\
&= \frac{1}{g'}\frac{\del g'}{\del y_1}u_{11}u_1 +\sum_{p=1}^{n}\frac{1}{g'}\frac{\del g'}{\del y_p}\frac{u_{p}}{u}
\end{aligned}
\]
On the other hand, by~\eqref{eq: LinearizedEquation} we have
\[
u_1\sum_p \frac{u_{1pp}}{u_{pp}}  + \frac{1}{g'}\frac{\del g'}{\del y_1}u_{11}u_1 = 0
\]
and so, in total
\[
u_1\sum_p \frac{u_{1pp}}{u_{pp}}-\sum_{p=1}^{n}\frac{1}{g'u_{11}}\frac{\del g'}{\del y_p}u_{11p}=\sum_{p=1}^{n}\frac{1}{g'}\frac{\del g'}{\del y_p}\frac{u_{p}}{u}.
\]
Substituting this expression into~\eqref{eq: PogorelovInequality1} yields
\begin{equation}
\begin{aligned}
0 &\geq \sum_{2 \leq p,q\leq n}\frac{u_{1pq}^2} {u_{pp}u_{qq}u_{11}}-u_{11}\frac{\del^2 \log g'}{\del y_1^2}  \\
  &\quad + \frac{n}{u} -\frac{u_1^2}{u^2u_{11}} + u_{11}+\sum_{p=1}^{n}\frac{1}{g'}\frac{\del g'}{\del y_p}\frac{u_{p}}{u}.
  \end{aligned}
\end{equation}
Since $u<0$, the structural property $(iv)$ gives
\[
\frac{\del^2 \log g'}{\del y_1^2} \leq 0 \quad \text{ and } \quad \frac{1}{g'}\sum_p\frac{\del g'}{\del y_p}\frac{u_{p}}{u} \geq \frac{\kappa}{u}.
\]
 Thus, in total, at $x_*$ we have
 \[
 0 \geq \frac{n+\kappa}{u} -\frac{u_1^2}{u^2u_{11}} + u_{11}.
 \]
 Multiplying through the $|u|^2u_{11}$ yields the result.

 We now address the case when $x_* \in \del\Omega \cap S_h \cap B_1(0)$.  First we remark that since $g,g'$ are smooth and bounded strictly away from zero, the result of Chen-Liu-Wang \cite{Chen-Liu-Wang}, and the regularity theory of Caffarelli \cite{Caffarelli2, Caffarelli3}, $u$ is smooth up to the boundary 
 
 By obliqueness, we can apply a shearing transformation in the $(x_2,\ldots,x_{n})$ directions which fixes the $x_n$ direction, we may assume that $e_n$ is the inward pointing normal vector to $\Omega$ at $x_*$.  In other words, we have $\frac{\del \phi_{\Omega}}{\del x_j}(x_*)=0$ for all $1 \leq j \leq n-1$.  Similarly, by a shearing transformation in the $(y_1,\ldots,y_n)$ variables which fixes the $y_n$ direction we may assume that  $e_n$ is the inward pointing normal vector to $\Omega'$ at $y_*=\nabla u(x_*)$.  In other words, $\frac{\del \phi_{\Omega'}}{\del y_j}(y_*)=0$ for all $1 \leq j \leq n-1$. 
 
 We may write the restriction of  $\Phi$ to $\del\Omega \cap S_h \cap B_1(0)$ as
 \[
 \Phi\big|_{\del \Omega}(x_1,x_2,\ldots,x_{n-1}) = \Phi(x_1,\ldots,x_{n-1},\phi_\Omega(x_2,\ldots,x_{n-1}))
 \]
 Since $\Phi$ is maximized at $x_*$ we have
\begin{equation} \label{eq: boundaryFirstDerivTest}
 0=\frac{\del \Phi}{\del x_j} + \frac{\del \Phi}{\del x_n}\frac{\del \phi_{\Omega}}{\del x_j} = \frac{\del \Phi}{\del x_j} \quad \text{ for }1\leq j\leq n-1.
 \end{equation}
 Differentiating a second time yields
 \begin{equation}\label{eq: boundarySecondDerivTest}
 \begin{aligned}
0&\geq  \frac{\del^2 \Phi}{\del x_j^2} + \frac{\del \Phi}{\del x_n}\frac{\del^2 \phi_{\Omega}}{\del x_j^2} + 2\frac{\del^2 \Phi}{\del x_n\del x_{j}}\frac{\del \phi_{\Omega}}{\del x_j} + \frac{\del ^2\Phi}{\del x_n^2}\left(\frac{\del \phi_{\Omega}}{\del x_j}\right)^2\\
&= \frac{\del^2 \Phi}{\del x_j^2} + \frac{\del \Phi}{\del x_n}\frac{\del^2 \phi_{\Omega}}{\del x_j^2}\bigg|_{x_*}
\end{aligned}
 \end{equation}
We now consider the derivative of $\Phi$ in the $x_n$ direction.  We have
 \[
 \Phi_{n} = \frac{u_{11n}}{u_{11}}+\frac{u_n}{u} + u_1u_{1n}
 \]
 On the other hand, from the boundary data we have
 \[
 u_{n} = \phi_{\Omega'}(u_1,\ldots,u_{n-1})
 \]
 and so
 \[
 \begin{aligned}
 u_{1n} &= \sum_{p=1}^{n-1}\frac{\del \phi_{\Omega'}}{\del y_p}u_{1p} \\
  u_{1n}(x_*) &= \sum_{p=1}^{n-1}\frac{\del \phi_{\Omega'}}{\del y_p}(y_*)u_{1p}(x_*) = 0.
  \end{aligned}
 \]
 Differentiating again yields
 \[
 u_{11n} = \sum_{1 \leq p,q \leq n-1}\frac{\del^2 \phi_{\Omega'}}{\del y_p\del y_q}u_{1p}u_{1q} + \sum_{p=1}^{n-1}\frac{\del \phi_{\Omega'}}{\del y_p}u_{11p}.
 \]
By~\eqref{eq: boundaryFirstDerivTest} we have
\[
\Phi_p = \frac{u_{11p}}{u_{11}} + \frac{u_p}{u} + u_{1p}u_1=0 \quad \text{ for } 1 \leq p \leq n-1
\]
and so
\[
\frac{u_{11n}}{u_{11}} = \sum_{1 \leq p,q \leq n-1}\frac{\del^2 \phi_{\Omega'}}{\del y_p\del y_q}\frac{u_{1p}u_{1q}}{u_{11}} - \sum_{p=1}^{n-1}\frac{\del \phi_{\Omega'}}{\del y_p}\left(\frac{u_p}{u} + u_{1p}u_1\right).
\]
This gives
\[
\begin{aligned}
\Phi_n &= \sum_{1 \leq p,q \leq n-1}\frac{\del^2 \phi_{\Omega'}}{\del y_p\del y_q}\frac{u_{1p}u_{1q}}{u_{11}} - \sum_{p=1}^{n-1}\frac{\del \phi_{\Omega'}}{\del y_p}\left(\frac{u_p}{u} + u_{1p}u_1\right) + \frac{\phi_{\Omega'}}{u} + u_1 \sum_{p=1}^{n-1}\frac{\del \phi_{\Omega'}}{\del y_p}u_{1p}\\
&= \sum_{1 \leq p,q \leq n-1}\frac{\del^2 \phi_{\Omega'}}{\del y_p\del y_q}\frac{u_{1p}u_{1q}}{u_{11}} +\frac{1}{|u|} \left(\sum_{p=1}^{n-1}\frac{\del \phi_{\Omega'}}{\del y_p}u_p - \phi_{\Omega'}\right).
\end{aligned}
\]
By the convexity of $\phi_{\Omega}$, the first term is non-negative. By the convexity of $\phi_{\Omega'}$ and $\phi_{\Omega'}(0) = 0$, the second term is non-negative as well. Hence we conclude that
\[
\Phi_{n}(x_*) \geq 0. 
\]
In particular, we must have
\begin{equation}\label{eq: normalDerivBoundaryPog}
\Phi_{n}(x_*)=0 \quad \text{ and } \quad  \Phi_{nn}(x_*) \leq 0.
\end{equation}
Substituting this into~\eqref{eq: boundarySecondDerivTest} we obtain
\begin{equation}\label{eq: nonNormalDerivBoundaryPog}
\frac{\del \Phi}{\del x_j}(x_*)=0 \quad \text{ and } \quad \frac{\del^2 \Phi}{\del x_j^2}(x_*) \leq 0. 
\end{equation}
Finally, since $u_{1n}(x_*)=0$, we may perform a shearing transformation preserving the $x_1$ and $x_n$ directions so that $D^2u(x_*)$ is diagonal.  By structural condition $(i)$, $\del \Omega'$ is invariant under translations in the $x_1$ direction and so equations~\eqref{eq: LinearizedEquation},~\eqref{eq: concavityEquation} hold up to the boundary.  Furthermore, by~\eqref{eq: normalDerivBoundaryPog} and ~\eqref{eq: nonNormalDerivBoundaryPog} we conclude that ~\eqref{eq: firstDerivTestPhi} and~\eqref{eq: secondDerivTestPhi} also hold at $x_*$.  We may therefore treat $x_*$ as if it were an interior point, and the previous argument leads to the desired estimate.
\end{proof}

\begin{thebibliography}{99}
\bibitem[3]{Caffarelli} Caffarelli, L. {\em The regularity of mappings with a convex potential.} J. Amer. Math. Soc. 5 (1992), no. 1, 99–104. 
\bibitem[4]{Caffarelli2} Caffarelli, L. {\em Boundary regularity of maps with convex potentials.} Comm. Pure Appl. Math. 45 (1992), no. 9, 1141--1151 
\bibitem[5]{Caffarelli3} Caffarelli, L. {\em Boundary regularity of maps with convex potentials. II} Ann. of Math. (2) 144 (1996), no. 3, 453–496. 
\bibitem[9]{Chen-Liu-Wang} Chen, S., Liu, J., X.-J. Wang {\em Global regularity for the Monge-Ampère equation with natural boundary condition.} Ann. of Math. (2) 194 (2021), no. 3, 745–793. 
\bibitem[18]{Jhaveri-Savin} Jhaveri, Y., Savin, O. {\em On the regularity of optimal transports between degenerate densities}, Arch. Ration. Mech. Anal. 245 (2022), no. 2, 819--861 
\bibitem[26]{Savin-Yu} Savin, O. Yu, H. {\em Regularity of optimal transport between planar convex domains. } Duke Math. J. 169 (2020), no. 7, 1305–1327. 
\bibitem[35]{Yuan} Yuan, Y. {\em A monotonicity approach to Pogorelov's Hessian estimates for Monge-Ampere equation. } Math. Eng. 5 (2023), no. 2, Paper No. 037, 6 pp.
\end{thebibliography}
\end{document}
